\clearpage\thispagestyle{empty}
\pdfbookmark[0]{Quellen zu Gauß’ Beobachtungen der Iris (September 1847)}{iris1847-sources}
\begin{center}\textbf{Quellen zu Gauß’ Beobachtungen der Iris (September 1847)}\end{center}
\noindent Diese KI-Übertragung ins heutige Deutsch folgt dem Bericht von Carl Friedrich Gauß auf S.~473 der \href{https://gdz.sub.uni-goettingen.de/id/PPN236007467}{\textit{Werke}, Band~VI} (Göttingen, 1874). Der Bericht ist auf den 7. September 1847 datiert. Der Erstdruck steht in den \textit{Nachrichten von der Georg-Augusts-Universität und der Königl. Gesellschaft der Wissenschaften zu Göttingen}, Nr.~11 vom \href{https://images.p.sub.uni-goettingen.de/iiif/image/gdz:PPN668664029_0003:00000163/full/1800,/0/default.jpg}{13. September 1847}, \href{https://images.p.sub.uni-goettingen.de/iiif/image/gdz:PPN668664029_0003:00000168/full/1800,/0/default.jpg}{S.~166}, \href{https://images.p.sub.uni-goettingen.de/iiif/image/gdz:PPN668664029_0003:00000169/full/1800,/0/default.jpg}{167} und \href{https://images.p.sub.uni-goettingen.de/iiif/image/gdz:PPN668664029_0003:00000170/full/1800,/0/default.jpg}{168}. Der Bericht beginnt auf S.~166 unter der Überschrift der Königlichen Gesellschaft der Wissenschaften.
\par\smallskip\noindent John Russell Hind entdeckte Iris. Gauß berichtet über seine eigenen Beobachtungen und unterscheidet Goldschmidts Elemente auf Grundlage der Beobachtungen von Gauß und Challis von der Bahn, die Gould eigens berechnete. Tabellenschreibweise und Zeichensetzung folgen hier der Werkausgabe. KI-Übertragung und KI-Satz.
\clearpage
